\documentclass[11pt]{article}
\usepackage{amsmath,amssymb,amsthm}
\newtheorem{theorem}{Theorem}
\title{A uniform relative-margin limit for mixtures of comparison orders}
\author{}\date{}
\begin{document}\maketitle

\paragraph{Status and scope.}
This is a complete written probabilistic proof, with exact finite
enumeration for the two smallest nontrivial square grids. It is not
Lean-checked. The theorem concerns weighted mixtures of order-comparison
signs, even allowing arbitrarily many votes. It does not bound arbitrary
affine-ratio sums and does not resolve either global HeadComplexity target.
Its purpose is to close a possible margin-based repair of the audited
per-vote reversal route.

For a strict total order on $r+s$ named objects
$R_1,\ldots,R_r,C_1,\ldots,C_s$, define the comparison table
$v\in\{-1,1\}^{r\times s}$ by $v_{ij}=1$ exactly when $R_i$ precedes
$C_j$. Let $\mathcal V_{r,s}$ be the set of all such tables. It contains
both constant tables and is closed under negation. Write $N=rs$ and
$M=|\mathcal V_{r,s}|\leq(r+s)!$.

\begin{theorem}
For every $r,s\geq1$ there is a sign table $\sigma$ such that every
finite weighted order mixture
\[
 S_{ij}=c+\sum_{h=1}^L\lambda_hv^{(h)}_{ij},
 \qquad v^{(h)}\in\mathcal V_{r,s},
\]
with $A=|c|+\sum_h|\lambda_h|>0$ obeys
\[
 \min_{i,j}\frac{\sigma_{ij}S_{ij}}A
 <\sqrt{\frac{2(\log M+1)}{rs}}
 \leq\sqrt{\frac{2(\log((r+s)!)+1)}{rs}}.                 \tag{1}
\]
All logarithms are natural. In particular, when $r=s=m\geq2$, the
right side is at most
\[
 \sqrt{\frac{4\log(2m)}m+\frac2{m^2}}
 =O\!\left(\sqrt{\frac{\log m}m}\right).                 \tag{2}
\]
\end{theorem}

\begin{proof}
Choose all $N$ entries of $\sigma$ independently and uniformly from
$\{-1,1\}$. For a fixed $v$, put $X_v=\sum_{ij}\sigma_{ij}v_{ij}$.
For every $u>0$, independence and
$\cosh u\leq\exp(u^2/2)$ give
\[
 \Pr(X_v\geq t)
 \leq\exp(-ut)\,\mathbb E\exp(uX_v)
 \leq\exp(-ut+Nu^2/2).
\]
Taking $u=t/N$ proves
$\Pr(X_v\geq t)\leq\exp(-t^2/(2N))$.
For $t=\sqrt{2N(\log M+1)}$, the union bound over the $M$ comparison
tables is at most $e^{-1}<1$. Thus some $\sigma$ satisfies
\[
             \sum_{ij}\sigma_{ij}v_{ij}<t
             \quad\hbox{for every }v\in\mathcal V_{r,s}.       \tag{3}
\]
For completeness, the elementary exponential inequality used above
follows by differentiating $u^2/2-\log\cosh u$: its derivative is
$u-\tanh u\geq0$ for $u\geq0$, because the derivative of
$u-\tanh u$ is $\tanh^2u\geq0$.

Absorb each coefficient sign into its order table by reversing that
order when necessary. Represent the constant term using the appropriate
constant comparison table. Applying (3) and adding with the resulting
nonnegative coefficients yields
\[
       \sum_{ij}\sigma_{ij}S_{ij}<tA.
\]
The strict inequality persists because at least one coefficient is
positive. The minimum of the $N$ numbers on the left after division
by $A$ is at most their average, proving (1). Finally
$\log((2m)!)\leq2m\log(2m)$ proves (2).
\end{proof}

\paragraph{Consequence for numerical vote approximation.}
Suppose a proposed universal order-mixture construction promised a
uniform relative margin bounded below by $1/(\log m)^K$, for some
fixed $K>0$. The theorem rules out that promise for all sufficiently
large square tables, regardless of the number of order votes or their
real weights. Indeed
$\sqrt{\log m/m}=o((\log m)^{-K})$.
Therefore replacing the known unit absolute majority margin by a
universal inverse-polylogarithmic relative margin cannot justify
coarse numerical reversal of the individual votes.

This conclusion does not rule out target-dependent signed cancellations,
joint reversal identities, nonuniform per-entry error estimates, or
other constructions using actual affine ratios. Nor does it prove
that the existing quantitative per-vote lower bound is sharp for every
achievable margin.

\paragraph{Exact finite check.}
The accompanying script enumerates all total orders and all target
tables for $m=2,3$. It computes
$\min_\sigma\max_{v\in\mathcal V_{m,m}}\sum_{ij}\sigma_{ij}v_{ij}$
using only integer bit counts, verifies negation and constant-table
closure, and checks the exact binomial correlation distribution for a
fixed order. The resulting maximum-correlation bounds for a worst
target are $2/4$ and $5/9$, respectively. These small enumerations
illustrate (3); the all-size
existence theorem is the probabilistic proof above.
\end{document}
